\documentclass[12pt, a4paper]{article}
\usepackage[utf8]{inputenc}
\usepackage[ngerman]{babel}
\usepackage{amsmath, amssymb}
\usepackage{geometry}
\geometry{margin=2.5cm}
\usepackage{parskip}

\title{\textbf{Aufgabe 11 -- Lösung} \\ Methoden I: Mathematik, 2026S}
\author{}
\date{}

\begin{document}
\maketitle

\section*{Angabe}

Kostenfunktion für die Herstellung von $Q$ Einheiten:
\[
    C(Q) = 12000 + 7Q + 0{,}0015\,Q^2 + 0{,}2\,Q^3
\]

% ============================================================
\section*{Teil 1: Kosten für 15 Einheiten}

Einfach $Q = 15$ einsetzen:

\begin{align*}
    C(15) &= 12000 + 7 \cdot 15 + 0{,}0015 \cdot 15^2 + 0{,}2 \cdot 15^3 \\[6pt]
    &= 12000 + 105 + 0{,}0015 \cdot 225 + 0{,}2 \cdot 3375 \\[6pt]
    &= 12000 + 105 + 0{,}3375 + 675 \\[6pt]
    &= \boxed{12780{,}3375}
\end{align*}

Die Herstellung von 15 Einheiten kostet \textbf{12.780,34} (Geldeinheiten).

% ============================================================
\section*{Teil 2: Kosten für eine zusätzliche Einheit (Grenzkosten)}

\textbf{Frage:} Wie viel kostet es, von $Q$ auf $Q + 1$ Einheiten zu gehen?

Das sind die \textbf{Mehrkosten}: $C(Q + 1) - C(Q)$.

\bigskip

\textbf{Schritt für Schritt:}

\begin{align*}
    C(Q+1) &= 12000 + 7(Q+1) + 0{,}0015(Q+1)^2 + 0{,}2(Q+1)^3
\end{align*}

Wir berechnen die Differenz $C(Q+1) - C(Q)$ für jeden Term einzeln:

\bigskip

\textbf{Konstante:} $12000 - 12000 = 0$

\textbf{Linearer Term:}
\[
    7(Q+1) - 7Q = 7
\]

\textbf{Quadratischer Term:}
\begin{align*}
    0{,}0015(Q+1)^2 - 0{,}0015\,Q^2 &= 0{,}0015\left[(Q+1)^2 - Q^2\right] \\
    &= 0{,}0015\left[Q^2 + 2Q + 1 - Q^2\right] \\
    &= 0{,}0015(2Q + 1)
\end{align*}

\textbf{Kubischer Term:}
\begin{align*}
    0{,}2(Q+1)^3 - 0{,}2\,Q^3 &= 0{,}2\left[(Q+1)^3 - Q^3\right]
\end{align*}

Nebenrechnung: $(Q+1)^3 = Q^3 + 3Q^2 + 3Q + 1$

\[
    (Q+1)^3 - Q^3 = 3Q^2 + 3Q + 1
\]

Also: $0{,}2(3Q^2 + 3Q + 1)$

\bigskip

\textbf{Alles zusammen:}

\begin{align*}
    C(Q+1) - C(Q) &= 7 + 0{,}0015(2Q + 1) + 0{,}2(3Q^2 + 3Q + 1)
\end{align*}

Ausmultiplizieren:

\begin{align*}
    &= 7 + 0{,}003Q + 0{,}0015 + 0{,}6Q^2 + 0{,}6Q + 0{,}2 \\[6pt]
    &= 0{,}6Q^2 + 0{,}603Q + 7{,}2015
\end{align*}

\[
    \boxed{C(Q+1) - C(Q) = 0{,}6Q^2 + 0{,}603Q + 7{,}2015}
\]

\bigskip

\textbf{Beispiel:} Zusätzliche Kosten von Einheit 15 auf 16:
\begin{align*}
    C(16) - C(15) &= 0{,}6 \cdot 225 + 0{,}603 \cdot 15 + 7{,}2015 \\
    &= 135 + 9{,}045 + 7{,}2015 = \boxed{151{,}2465}
\end{align*}

\textit{Die 16. Einheit kostet also ca. 151,25 mehr als die 15. Es wird immer teurer, weil der kubische Term dominiert.}

\end{document}
